%Data institusi

%\input{Dozentdaten}

%\renewcommand{\fachbereich}{Bidang studi}

%\renewcommand{\dozent}{Prof. Dr. . }

%Data ujian

\renewcommand{\klausurgebiet}{ }

\renewcommand{\klausurtyp}{ }

\renewcommand{\klausurdatum}{. 20}

\klausurvorspann {\fachbereich} {\klausurdatum} {\dozent} {\klausurgebiet} {\klausurtyp}

%Data untuk tabel poin berikut

\renewcommand{\aeins}{  3  }

\renewcommand{\azwei}{  3   }

\renewcommand{\adrei}{  6  }

\renewcommand{\avier}{  5  }

\renewcommand{\afuenf}{  5  }

\renewcommand{\asechs}{  2  }

\renewcommand{\asieben}{  5  }

\renewcommand{\aacht}{  5  }

\renewcommand{\aneun}{  7  }

\renewcommand{\azehn}{  3  }

\renewcommand{\aelf}{  2  }

\renewcommand{\azwoelf}{  5  }

\renewcommand{\adreizehn}{  9   }

\renewcommand{\avierzehn}{  5  }

\renewcommand{\afuenfzehn}{  65  }

\renewcommand{\asechzehn}{    }

\renewcommand{\asiebzehn}{    }

\renewcommand{\aachtzehn}{    }

\renewcommand{\aneunzehn}{    }

\renewcommand{\azwanzig}{    }

\renewcommand{\aeinundzwanzig}{    }

\renewcommand{\azweiundzwanzig}{    }

\renewcommand{\adreiundzwanzig}{    }

\renewcommand{\avierundzwanzig}{    }

\renewcommand{\afuenfundzwanzig}{    }

\renewcommand{\asechsundzwanzig}{    }

\punktetabellevierzehn

\klausurnote

\newpage

\setcounter{section}{0}

\inputaufgabegibtloesung
{3}
{

Definisikan istilah-istilah berikut
\zusatzklammer {yang dicetak miring} {} {}.
\aufzaehlungsechs{\stichwort {Lingkaran kelengkungan} {}
dari suatu
\definitionsverweis {kurva}{}{}
yang dua kali
\definitionsverweis {terdiferensialkan secara kontinu}{}{}
dan
\definitionsverweis {berparametrisasi panjang busur}{}{}.
Misalkan
\maabbdisp {\gamma} {I} {\R^2
} {}
dan ambil titik

\mavergleichskette
{\vergleichskette
{ t_0
}
{ \in }{ I
}
{  }{
}
{    }{
}
{   }{
}
}
{}{}{.}

}{Suatu \stichwort {submanifold tertutup} {}
\mathl{M \subseteq N}{} dari manifold diferensiabel $N$.

}{\stichwort {Ruang singgung} {} pada titik
\mathl{P\in M}{} dari manifold diferensiabel $M$.

}{Istilah
\stichwort {berorientasi} {}
untuk suatu
\definitionsverweis {manifold diferensiabel}{}{}
$M$.

}{\stichwort {Batas} {}
suatu
\definitionsverweis {manifold berbatas}{}{.}

}{\stichwort {Simbol Christoffel} {}

\mathl{\Gamma_{ij}^k}{} bagi koneksi Levi-Civita dari struktur Riemann $\left(  g_{ij}  \right)$ pada himpunan terbuka

\mavergleichskette
{\vergleichskette
{ U
}
{ \subseteq }{ \R^n
}
{  }{
}
{    }{
}
{   }{
}
}
{}{}{.}
}

}
{} {}

\inputaufgabegibtloesung
{3}
{

Nyatakan teorema-teorema berikut.
\aufzaehlungdrei{Teorema tentang citra
\definitionsverweis {pemetaan Gauss}{}{}
dari suatu hipermuka.}{Rumus luas permukaan dari permukaan putar yang dibentuk oleh
\definitionsverweis {kurva diferensiabel}{}{}
\maabbeledisp {\gamma} {]a,b[} {\R^2
} { t } {(x(t),y(t))
} {,}
dengan

\mavergleichskette
{\vergleichskette
{  y(t)
}
{ > }{  0
}
{  }{
}
{    }{
}
{   }{
}
}
{}{}{.}}{\stichwort {Aturan hasil kali} {}
untuk turunan eksterior.}

}
{} {}

\inputaufgabegibtloesung
{6 (1+2+3)}
{

Misalkan $Y$ adalah hipermuka yang diberikan oleh persamaan

\mavergleichskettedisp
{\vergleichskette
{ 2x^2 +3y^2+5z^2
}
{ =} { 10
}
{ } {
}
{ } {
}
{ } {
}
}
{}{}{}
di $\R^3$. Misalkan

\mavergleichskette
{\vergleichskette
{ P
}
{ = }{ \begin{pmatrix}  -1 \\ -1\\  1   \end{pmatrix}
}
{ \in }{ Y
}
{    }{
}
{   }{
}
}
{}{}{}
dan

\mavergleichskette
{\vergleichskette
{ w
}
{ = }{ \begin{pmatrix}  -3 \\ 2 \\  4   \end{pmatrix}
}
{ \in }{ \R^3
}
{    }{
}
{   }{
}
}
{}{}{.}
\aufzaehlungdrei{Tentukan
\definitionsverweis {ruang singgung}{}{}

\mathl{T_PY}{.}
}{Tentukan
\definitionsverweis {dekomposisi ortogonal}{}{}
$w$ menjadi komponen tangensial dan komponen ortogonal pada titik $P$.
}{Realisasikan komponen tangensial $w$ di $T_PY$ melalui suatu kurva diferensiabel yang seluruhnya terletak pada $Y$.
}

}
{} {}

\inputaufgabegibtloesung
{5}
{

Misalkan
\maabbeledisp {\gamma} {\R} { \R^2
} { t } { \begin{pmatrix}  \cos \left(  \omega t \right)  \\ \sin  \left(  \omega t \right)    \end{pmatrix}
} {,}
dengan

\mavergleichskette
{\vergleichskette
{ \omega
}
{ > }{ 0
}
{  }{
}
{    }{
}
{   }{
}
}
{}{}{.}
Tunjukkan bahwa tidak ada gerak melingkar lain dengan norma kecepatan konstan yang memiliki nilai dan turunan hingga orde kedua yang sama dengan $\gamma$ di titik parameter

\mavergleichskette
{\vergleichskette
{ t
}
{ = }{ 0
}
{  }{
}
{    }{
}
{   }{
}
}
{}{}{}
tersebut.

}
{} {}

\inputaufgabegibtloesung
{5}
{

Buktikan teorema tentang kelengkungan suatu kurva bidang yang diberikan secara implisit.

}
{} {}

\inputaufgabegibtloesung
{2}
{

Berikan contoh manifold diferensiabel
\mathkor {} {M} {dan} {N} {}
dengan

\mavergleichskette
{\vergleichskette
{ \operatorname{ dim}_{  } ^{  }  \left(   M   \right), \, \operatorname{ dim}_{  } ^{  }  \left(   N   \right)
}
{ \geq }{ 1
}
{  }{
}
{    }{
}
{   }{
}
}
{}{}{}
serta pemetaan diferensiabel
\maabb {\varphi} { M } { N
} {}
dan
\maabb {\psi} { N } { M
} {}
sedemikian sehingga

\mavergleichskette
{\vergleichskette
{ \psi \circ \varphi
}
{ = }{
\operatorname{Id}_{ M }
}
{  }{
}
{    }{
}
{   }{
}
}
{}{}{}
dan

\mavergleichskette
{\vergleichskette
{\varphi \circ  \psi
}
{ \neq }{
\operatorname{Id}_{ N }
}
{  }{
}
{    }{
}
{   }{
}
}
{}{}{}
berlaku.

}
{} {}

\inputaufgabegibtloesung
{5}
{

Buktikan bahwa ekuivalensi tangensial lintasan-lintasan pada suatu manifold di titik $P$ dapat diperiksa dengan menggunakan sebarang peta koordinat.

}
{} {}

\inputaufgabegibtloesung
{5}
{

Tunjukkan bahwa pemetaan singgung $T(\varphi)$ dari
\maabbeledisp {\varphi} {\R^1} {S^1
} { t } {( \cos  t, \sin  t  )
} {,}
bersifat surjektif.

}
{} {}

\inputaufgabegibtloesung
{7 (3+4)}
{

Tinjau pemetaan diferensiabel
\maabbeledisp {\gamma} { [1,c] } {\R^2
} { t } {(t,t^{3})
} {,}
\zusatzklammer {dengan \mathlk{c \geq 1}{}} {} {}
dan
\maabbeledisp {\varphi} {\R_+ \times \R_+ } {\R^3
} {(u,v)} {(u^3,u^2+v^2,u^{-1}v^{-1} )
} {,}
serta bentuk diferensial

\mavergleichskettedisp
{\vergleichskette
{ \omega
}
{ =} { (x-y)dx-z^2dy+dz
}
{ } {
}
{ } {
}
{ } {
}
}
{}{}{}
pada $\R^3$.

a) Hitung bentuk diferensial tarik balik
\mathl{\varphi^* \omega}{} pada $\R_+ \times \R_+$.

b) Hitung integral lintasan dari bentuk diferensial
\mathl{\varphi^* \omega}{} sepanjang lintasan $\gamma$ sebagai fungsi dari $c$.

}
{} {}

\inputaufgabegibtloesung
{3}
{

Misalkan
\maabbdisp {\psi} { L } { M
} {}
adalah
\definitionsverweis {pemetaan diferensiabel}{}{}
antara
\definitionsverweis {manifold diferensiabel}{}{.}
Selanjutnya, misalkan
\maabbdisp {f} { M } {\R
} {}
adalah fungsi pada $M$, dan $\omega$ adalah
$k$-\definitionsverweis {bentuk}{}{}
pada $M$. Tunjukkan bahwa

\mavergleichskettedisp
{\vergleichskette
{ \psi^*(f \omega)
}
{ =} { \psi^*(f)  \psi^* \omega
}
{ } {
}
{ } {
}
{ } {
}
}
{}{}{.}

}
{} {}

\inputaufgabe
{2}
{

Deskripsikan sebanyak mungkin
\definitionsverweis {manifold berbatas}{}{}
berdimensi satu yang
\definitionsverweis {terhubung}{}{.}

}
{} {}

\inputaufgabegibtloesung
{5}
{

Misalkan
\maabbdisp {f} { [a,b] } { \R_{+}
} {}
adalah
\definitionsverweis {fungsi yang terdiferensialkan secara kontinu}{}{.}
Kita pandang
\definitionsverweis {subgraf}{}{}
sebagai
\definitionsverweis {manifold berbatas}{}{}.
Batasnya terdiri dari grafik, interval dasar, dan kedua sisi tegak, tetapi tidak mencakup keempat titik sudut. Buktikan secara langsung berlakunya teorema Stokes untuk
\definitionsverweis {bentuk diferensial}{}{}

\mavergleichskette
{\vergleichskette
{ \omega
}
{ = }{ xdy
}
{  }{
}
{    }{
}
{   }{
}
}
{}{}{.}

}
{} {}

\inputaufgabegibtloesung
{9}
{

Buktikan teorema titik tetap Brouwer.

}
{} {}

\inputaufgabegibtloesung
{5 (3+1+1)}
{

Tinjau silinder

\mavergleichskette
{\vergleichskette
{ Z
}
{ = }{ S^1 \times \R
}
{ \subseteq }{ \R^2 \times \R
}
{    }{
}
{   }{
}
}
{}{}{}
dengan
\definitionsverweis {koneksi trivial}{}{}
pada
\definitionsverweis {bundel tangen}{}{.}
\aufzaehlungdrei{Tentukan, dalam suatu peta isometrik yang sesuai,
\definitionsverweis {persamaan diferensial geodesik}{}{}
untuk kurva geodesik $\psi$ pada $M$ yang memenuhi

\mavergleichskettedisp
{\vergleichskette
{ \psi(0)
}
{ =} { (1,0,3)
}
{ } {
}
{ } {
}
{ } {
}
}
{}{}{}
dan

\mavergleichskettedisp
{\vergleichskette
{ \psi'(0)
}
{ =} { (0,1,-4)
}
{ } {
}
{ } {
}
{ } {
}
}
{}{}{}
.
}{Selesaikan persamaan diferensial dari (1) dalam peta tersebut.
}{Carilah suatu kurva geodesik di $Z$ yang memenuhi syarat-syarat pada (1).
}

}
{} {}
